\documentclass[a4paper,12pt]{article}
\marginparwidth 0 pt
\oddsidemargin 0 truecm
\evensidemargin  0 pt
\marginparsep 0 pt
\topmargin 0 truecm
\headsep 0 pt
\textheight 23 truecm
\textwidth 14.9 truecm
\hoffset 0.5 truecm
\footskip 60 pt
\usepackage{fancyvrb}%\begin{Verbatim}[frame=single]
\usepackage{amsmath,enumitem}
\usepackage{amssymb}
\usepackage{amsfonts}
\usepackage{latexsym}
\usepackage{graphicx}
\usepackage{fancyhdr}
\setlength{\headheight}{15.2pt}
\renewcommand{\headrulewidth}{0pt} % Remove head line.
\pagestyle{fancy}
\usepackage{lastpage}
\cfoot{P.\thepage\ of \pageref{LastPage} }
\usepackage{xspace}
%%\usepackage{layout}

\def \la{{\langle}}
\def \ra{{\rangle}}
\def \R{{\mathbb  R}}
\def \P{{\mathbb  P}}
\def \C{{\mathbb  C}}
\def \H{{\mathbb  H}}
\def \Z{{\mathbb  Z}}
\newcommand{\N}{{\Bbb N}}
%\def \displaystyle\int{{\displaystyle\int}}

\def \i{{\bf i}}
\def \fd{{\bullet}}
\def \proof{{\noindent {\bf Proof.} \ \ }}
\newcommand{\qed}{\begin{flushright} {\bf Q.E.D.}\ \ \ \ \
                  \end{flushright} }

\newcommand{\MySoln}[1]{ {#1}\xspace }
\renewcommand{\MySoln}[1]{ \xspace }  %Remove2ShowQuestion-1/3%
\newcommand{\MySn}{\newpage $\ $}
\renewcommand{\MySn}{ \xspace }  %Remove2ShowQuestion-2/3%
%Comment out below for solution: \baselineskip=22pt  %Remove2ShowQuestion-3/3%

\begin{document}

\centerline{\bf THE UNIVERSITY OF HONG KONG}
\bigskip
\centerline{\bf DEPARTMENT OF MATHEMATICS}

\vskip 2truecm
\noindent
MATH4907: NUMERICAL METHODS FOR FINANCIAL CALCULUS

\vskip 1truecm
\noindent
20 October, 2026 \hfill 1:05p.m. -- 2:40p.m.

\it
\vskip 1truecm
\noindent
Only approved calculators as announced by the Examinations Secretary can be
used in this examination.  It is candidates' responsibility to ensure that
their calculator operate satisfactorily, and candidates must record the name
and type of the calculator used on the front page of the examination script.


\vskip 2truecm

\rm
\centerline{Answer ALL {\bf 4} questions}

\vskip 1truecm

\noindent

{\bf Read each question carefully. You may not be credited any score if you put down your conclusion ONLY.}

\hfill\\

{\bf Instructions}

\begin{itemize}%[label=(\roman*)]%[(i)]
%\item It is an OPEN BOOK test but you MUST COMPLETE THE WORK BY YOURSELVES.
%\item Please hand write your solutions in blank paper. Typed solution would not be accepted.
%\item Please include the calculations in your solution neatly. This requirement is to ensure that you are not using any mathematics software (which is forbidden).
%\item Please follow the exam policy by HKU and place two cameras from the side and front. Record your screen, and send to me by email after the exam. 
\item You should stop working at 2:40pm.
%\item Prepare and upload a soft copy of your solutions from 5:00pm to 5:10pm. You should prepare a single pdf file (less than 10MB) of solutions and {\bf upload your PDF file via the designated MOODLE drop box before 5:10pm}.
%\item {\it You are still allowed to upload the pdf files in the grace period 1545-1615}. However, you must write to the teacher to explain your late submission, with the completed ``Late Submission Investigation Form'' before the deadline (would be announced after the test). {\it If the excuse is not accepted, the candidate would be disqualified in this test.}
\item {\bf No submission would be accepted after 2:40pm}. 
\end{itemize}


{\bf Note:} You should always give precise and adequate explanations to support your
conclusions.  Clarity of presentation of your argument counts.  So {\bf think carefully before you
write.}

\newpage
\baselineskip=22pt  %Remove2ShowQuestion-3/3%


\bigskip

\begin{enumerate} 
\item (20 points) For each question, select \textbf{all} correct answers.  
There may be more than one correct answer for each question.
\begin{enumerate}[label=(\arabic*)]
\item Which statements about arbitrage and the no-arbitrage principle are correct?
\begin{enumerate}[label=\Alph*.]
    \item An arbitrage strategy may have non-positive initial value.
    \item An arbitrage strategy has zero probability of loss at some future time.
    \item An arbitrage strategy must guarantee strictly positive profit in every state.
    \item If two portfolios have the same terminal payoff almost surely, then they should have the same initial value under no-arbitrage.
    \item No-arbitrage is one of the fundamental principles used to price contingent claims.
\end{enumerate}

\textbf{Correct answers: A, B, D, E}
\item In the Black--Scholes model, which statements are correct?

\begin{enumerate}[label=\Alph*.]
    \item The underlying asset price follows
    \[
    dS_t=\mu S_t\,dt+\sigma S_t\,dW_t.
    \]
    \item \(W_t\) denotes a Wiener process.
    \item The derivative value is written as \(V_t=V(t,S_t)\).
    \item The volatility parameter is usually denoted by \(r\).
    \item The risk-free bond grows as
    \[
    B_t=B_0e^{rt}.
    \]
\end{enumerate}

\textbf{Correct answers: A, B, C, E}
\item Which statements about the binomial tree model are correct?

\begin{enumerate}[label=\Alph*.]
    \item In one time step, the stock price can move to either \(uS_i\) or \(dS_i\).
    \item The no-arbitrage condition is
    \[
    d<e^{r\Delta t}<u.
    \]
    \item The risk-neutral probability is
    \[
    q=\frac{e^{r\Delta t}-d}{u-d}.
    \]
    \item The real-world probability \(p\) must be used directly for derivative pricing.
    \item The derivative value can be computed by backward induction.
\end{enumerate}

\textbf{Correct answers: A, B, C, E}
\item Which statements about the trinomial tree model are correct?

\begin{enumerate}[label=\Alph*.]
    \item The stock price may move up, down, or remain unchanged.
    \item In log-price form, the movement may be written as
    \[
    x_{i+1}\in\{x_i+\Delta x,x_i,x_i-\Delta x\}.
    \]
    \item The probabilities are often denoted by \(p_u,p_0,p_d\).
    \item The probabilities must satisfy
    \[
    p_u+p_0+p_d=1.
    \]
    \item The trinomial tree is stable for every choice of \(\Delta x\).
\end{enumerate}

\textbf{Correct answers: A, B, C, D}
\item Which finite difference formulas are correct?

\begin{enumerate}[label=\Alph*.]
    \item Forward difference:
    \[
    f'(x)\approx \frac{f(x+h)-f(x)}{h}.
    \]
    \item Central difference for the first derivative:
    \[
    f'(x)\approx \frac{f(x+h)-f(x-h)}{2h}.
    \]
    \item Central difference for the second derivative:
    \[
    f''(x)\approx \frac{f(x-h)-2f(x)+f(x+h)}{h^2}.
    \]
    \item The forward difference approximation has error \(O(h^2)\).
    \item The central difference approximation for \(f'(x)\) has error \(O(h^2)\).
\end{enumerate}

\textbf{Correct answers: A, B, C, E}
\item For the heat equation
\[
u_\tau=u_{xx},
\]
with
\[
\lambda=\frac{k}{h^2},
\]
which statements about the explicit finite difference scheme are correct?

\begin{enumerate}[label=\Alph*.]
    \item The explicit scheme computes the next time level directly from the current time level.
    \item The explicit scheme can be written as
    \[
    u^{j+1}=(I-\lambda B)u^j+\lambda c^j.
    \]
    \item The explicit scheme requires solving a matrix equation at every time step.
    \item The usual stability condition is
    \[
    0<\lambda\leq \frac12.
    \]
    \item If \(\lambda\) is too large, the numerical solution may become unstable.
\end{enumerate}

\textbf{Correct answers: A, B, D, E}
\item Which statements about the implicit finite difference scheme are correct?

\begin{enumerate}[label=\Alph*.]
    \item It may require solving a linear system at each time step.
    \item For the heat equation, one form is
    \[
    (I+\lambda B)u^j=u^{j-1}+\lambda c^j.
    \]
    \item It is called implicit because unknown future time-level values appear in the equation.
    \item The matrix \(I+\lambda B\) is strictly diagonally dominant for \(\lambda>0\).
    \item The implicit scheme is always identical to the explicit scheme.
\end{enumerate}

\textbf{Correct answers: A, B, C, D}
\item Which statements about consistency, stability, and convergence are correct?

\begin{enumerate}[label=\Alph*.]
    \item Consistency means the finite difference scheme approximates the PDE locally as the grid is refined.
    \item Stability means numerical errors do not grow uncontrollably during computation.
    \item Convergence means the numerical solution approaches the exact solution as \(\Delta x,\Delta t\to0\).
    \item The Lax equivalence theorem relates consistency, stability, and convergence for suitable linear problems.
    \item A consistent scheme is always convergent, even if it is unstable.
\end{enumerate}

\textbf{Correct answers: A, B, C, D}
\item Which statements about lattice methods and finite difference methods are correct?

\begin{enumerate}[label=\Alph*.]
    \item Tree models can be interpreted as particular explicit finite difference schemes.
    \item A trinomial tree may be obtained from a finite difference discretization of the transformed Black--Scholes equation.
    \item If
    \[
    h^2=k\sigma^2,
    \]
    then \(p_0=0\), giving a binomial tree model.
    \item The probabilities \(p_u,p_0,p_d\) must be nonnegative to be interpreted as probabilities.
    \item Lattice methods and finite difference methods are completely unrelated.
\end{enumerate}

\textbf{Correct answers: A, B, C, D}
\item Which statements about Fourier stability analysis are correct?

\begin{enumerate}[label=\Alph*.]
    \item Fourier analysis can be used to study stability of finite difference schemes with constant coefficients.
    \item The amplification factor describes how Fourier modes evolve from one time level to the next.
    \item For the explicit heat scheme, the amplification factor is
    \[
    g(\xi)=1-4\lambda\sin^2\left(\frac{\xi}{2}\right).
    \]
    \item Stability requires that high-frequency errors do not grow uncontrollably.
    \item The explicit heat scheme is stable for all \(\lambda>0\).
\end{enumerate}

\textbf{Correct answers: A, B, C, D}
\end{enumerate}
\newpage
\item (15 points) Define a European call option with strike price $K$, expiry $T$, and the price of underlying asset following the Geometric Brownian motion, 
\begin{align*}
\mathrm{d}S_t=\mu S_t\mathrm{d}t+\sigma S_t\mathrm{d}W_t.
\end{align*}
Let the risk neutral interest rate be $r$. Given the parameters
\[ S_0 = K = 100,\ r=0.05,\ \sigma=0.15,\ T=1. \]
{Use $2$--step CRR binomial tree to price a European call option. }
\begin{enumerate}
 \item (5 points) Compute the up factor $u$ and down factor $d$.
 \item (5 points) Compute the risk-neutral probability $p$.
  \item (5 points) Compute the value of the European call option.
\end{enumerate}
%\begin{solution}
\begin{proof}\mbox{}
We use a 2-step CRR binomial tree. Hence
\[
N=2,\qquad \Delta t=\frac{T}{N}=\frac{1}{2}.
\]

\begin{enumerate}
    \item In the CRR model, the up factor and down factor are
    \[
    u=e^{\sigma\sqrt{\Delta t}},
    \qquad
    d=e^{-\sigma\sqrt{\Delta t}}.
    \]
    Substituting the given parameters,
    \[
    u=e^{0.15\sqrt{0.5}}
    \approx 1.1119,
    \]
    and
    \[
    d=e^{-0.15\sqrt{0.5}}
    \approx 0.8994.
    \]
    Therefore,
    \[
    \boxed{u\approx 1.1119,\qquad d\approx 0.8994.}
    \]

    \item The risk-neutral probability is
    \[
    p=\frac{e^{r\Delta t}-d}{u-d}.
    \]
    Thus,
    \[
    p
    =
    \frac{e^{0.05(0.5)}-0.8994}{1.1119-0.8994}.
    \]
    Since
    \[
    e^{0.025}\approx 1.0253,
    \]
    we get
    \[
    p
    \approx
    \frac{1.0253-0.8994}{1.1119-0.8994}
    \approx 0.5926.
    \]
    Hence,
    \[
    \boxed{p\approx 0.5926.}
    \]

    \item At maturity, the possible stock prices are
    \[
    S_{uu}=S_0u^2,
    \qquad
    S_{ud}=S_0ud,
    \qquad
    S_{dd}=S_0d^2.
    \]
    Since \(S_0=100\), we have
    \[
    S_{uu}=100(1.1119)^2\approx 123.6311,
    \]
    \[
    S_{ud}=100(1.1119)(0.8994)\approx 100,
    \]
    and
    \[
    S_{dd}=100(0.8994)^2\approx 80.8858.
    \]

    The payoff of a European call option is
    \[
    C_T=(S_T-K)^+.
    \]
    Therefore, the terminal payoffs are
    \[
    C_{uu}=(123.6311-100)^+=23.6311,
    \]
    \[
    C_{ud}=(100-100)^+=0,
    \]
    and
    \[
    C_{dd}=(80.8858-100)^+=0.
    \]

    Using backward induction,
    \[
    C_0
    =
    e^{-rT}
    \left[
    p^2 C_{uu}
    +2p(1-p)C_{ud}
    +(1-p)^2C_{dd}
    \right].
    \]
    Since \(C_{ud}=C_{dd}=0\), this simplifies to
    \[
    C_0=e^{-0.05}p^2C_{uu}.
    \]
    Substituting the values,
    \[
    C_0
    =
    e^{-0.05}(0.5926)^2(23.6311).
    \]
    Hence,
    \[
    C_0\approx 7.8945.
    \]

    Therefore, the 2-step CRR binomial tree price of the European call option is
    \[
    \boxed{C_0\approx 7.89.}
    \]

\end{enumerate}

\end{proof}

\newpage
\item (20 points) Consider the initial boundary value problem (IBVP)
\[ \left\{ \begin{array}{ll}
u_t - 2u_{xx} + u = 0, &\quad \mbox{for all $(x,t)\in(0,1)\times(0,+\infty)$}\\
u(x,0) = \phi_0(x), &\quad \mbox{for all $x \in (0,1)$}\\
u(0,t)=u(1,t) = 0, \quad \mbox{for all $t > 0$}
\end{array}\right.. \]
for some function $\phi_0 : (0,1) \to \mathbb R$. 

It's decided to find a numerical solution to this IBVP over $[0,1]\times[0,T]$ for some fixed $T>0$ as following. Set $k=T/m$ and $h=1/(n+1)$ for some $m, n \in \mathbb N$; and set $x_i=ih$ and $t_j=jk$. Approximate $u_t(x_i,t_j)$ by a forward finite difference and approximate $u_{xx}(x_i,t_j)$ by the centered finite difference.

\begin{enumerate}
\item Put down the finite difference scheme in terms of $v_{i,j}$, the approximation of $u(x_i,t_j)$.
\item Find the amplification factor of this finite difference scheme. Hence, show that the scheme is  stable with respect to $\|\cdot\|_{2,\Delta x}$. If it is not unconditionally stable, please state clearly the condition.
\item 
\begin{enumerate}[label=(\roman*)]
\item Show that the scheme is consistent with the pde with respect to $\|\cdot\|_{\infty}$. State the assumptions you have made.
\item Express the finite difference scheme in matrix form. Hence, or otherwise, determine whether the scheme is convergent with respect to $\|\cdot\|_\infty$.
\end{enumerate}
\end{enumerate}


\begin{proof}\mbox{}
\begin{enumerate}
\item With the suggested approximation by finite differences, we have
\begin{gather*}
\frac{v_{i,j+1}-v_{i,j}}k - 2 \frac{v_{i-1,j}-2v_{i,j}+v_{i+1,j}}{h^2} + v_{i,j} = 0\\
v_{i,j+1}-v_{i,j} - 2\lambda \left( v_{i-1,j}-2v_{i,j}+v_{i+1,j} \right) + k v_{i,j} = 0\\
v_{i,j+1} = 2\lambda v_{i-1,j} + (1-k-4\lambda) v_{i,j} + 2\lambda v_{i+1,j}
\end{gather*}
where $\lambda=k/h^2$.

\item The finite difference scheme is equivalent to
\[ \mathbf v_{j+1} = \left( 2\lambda \tau_{-1} + (1-k-4\lambda) \tau_0 + 2\tau_{+1} \right) \mathbf v_j. \]
Apply the discrete Fourier transform to the above equation to get
\begin{eqnarray*}
\mathcal F(\mathbf v_{j+1})(\xi) &=& \left( 2\lambda e^{-\sqrt{-1}\xi} + (1-k-4\lambda) + 2\lambda e^{\sqrt{-1}\xi} \right)\mathcal F(\mathbf v_{j})(\xi)\\
&=& \left[ 1 - k - 4\lambda + 4\lambda\cos\xi \right] F(\mathbf v_{j})(\xi)\\
&=& \left[ 1 - k - 4\lambda \left( 1 - \cos\xi \right) \right] F(\mathbf v_{j})(\xi)\\
&=& \left( 1 - k - 8\lambda \sin^2\frac\xi 2 \right) F(\mathbf v_{j})(\xi)
\end{eqnarray*}
So the amplification factor is
\[ g(\xi) = \left( 1 - 8\lambda \sin^2\frac\xi 2 \right) - k. \]
Clearly, $\displaystyle 1 - 8\lambda \sin^2\frac\xi 2 \le 1$. On the other hand,
\[ 1 - 8\lambda \sin^2\frac\xi 2 \ge -1 \quad\iff\quad \lambda \le \frac 14. \]
Assume that $0 < \lambda \le 1/4$. Then
\[ |g(\xi)| \le \left| 1 - 8\lambda \sin^2\frac\xi 2 \right| + k \le 1+k \le e^{k}. \]
It follows that the scheme is stable with respect to $\|\cdot\|_{2,\Delta x}$ (with $K=\beta=1$).

\item 
\begin{enumerate}[label=(\roman*)]
\item Let $u$ be a solution of the pde and set $u_{i,j}=u(x_i,t_j)$. Apply the Taylor's theorem to get
\begin{eqnarray*}
\tau_{i,j} &=& \frac{u_{i,j+1}-u_{i,j}}k - 2 \frac{u_{i-1,j}-2u_{i,j}+u_{i+1,j}}{h^2} + u_{i,j}\\
&=& (u_t)_{i,j} + \frac 12 u_{tt}(x_i,(1-s_1)t_i+s_1t_{i+1}) t  + u_{i,j} \\
&&\ - 2 \left( 
(u_{xx})_{i,j} + \frac 1{24} \left( u_{xxxx}((1-s_2)x_{i-1}+s_2x_i) + u_{xxxx}((1-s_3)x_{i-1}+s_3x_i) \right) h^2 \right)\\
&=& (u_t-2u_{xx}+u)_{i,j} + \frac 12 u_{tt}(x_i,(1-s_1)t_i+s_1t_{i+1}) k\\
&&\ - \frac 1{12} \left( u_{xxxx}((1-s_2)x_{i-1}+s_2 x_i) + u_{xxxx}((1-s_3)x_{i-1}+s_3 x_i) \right) h^2 \\
&=& \frac 12 u_{tt}(x_i,(1-s_1)t_i+s_1t_{i+1}) k\\
&&\ - \frac 1{12} \left( u_{xxxx}((1-s_2)x_{i-1}+s_2 x_i) + u_{xxxx}((1-s_3)x_{i-1}+s_3 x_i) \right) h^2
\end{eqnarray*}
for some $0 < s_1,s_2,s_3 < 1$.

Let $R=[0,1]\times[0,T]$. Under the assumptions that $u_{tt}$ and $u_{xxxx}$ are bounded over $R$, we have
\begin{gather*}
|\tau_{i,j}| \le \frac 12 \max_{(x,t)\in R}\ |u_{tt}(x,t)|\ k + \frac 16 \max_{(x,t)\in R}\ |u_{xxxx}(x,t)|\ h^2\\
\|\boldsymbol\tau_j\|_\infty \le \frac 12 \max_{(x,t)\in R}\ |u_{tt}(x,t)|\ k + \frac 16 \max_{(x,t)\in R}\ |u_{xxxx}(x,t)|\ h^2.
\end{gather*}
It follows that $\|\boldsymbol\tau_j\|_\infty=O(h^2)+O(k)$ so that the scheme is consistent with the pde with respect to $\|\cdot\|_\infty$.

\item Using the boundary condition, $v_{0,j} = v_{n+1,j} = 0$ for all $j$. The finite difference scheme in matrix form is
\[ \mathbf v_{j+1} = A\mathbf x_j \]
where $A=(a_{ij})$ is a tridiagonal matrix with
\[ a_{ij} = \left\{ \begin{array}{ll}
1-k-4\lambda &\ \mbox{if $i=j$}\\
2\lambda &\ \mbox{if $|i-j|=1$}\\
0 &\ \mbox{otherwise}
\end{array}\right.. \]
Note that
\[ 1-k-4\lambda \ge 0 \ \iff\ 1-k\left(1-\frac 4{h^2}\right) \ge 0 \ \iff\ k \le \frac{h^2}{4+h^2} (<1). \]
It follows that if $0 < k \le \dfrac{h^2}{4+h^2}$, then
\[ \|A\|_\infty = \max\left\{ |1-k-4\lambda|+|2\lambda|,|1-k-4\lambda|+2|2\lambda| \right\} = 1-k-4\lambda+4\lambda = 1-k \in (0,1) \]
so that the scheme is stable with respect to $\|\cdot\|_\infty$ ($A$ is symmetric).

Lastly, as the scheme is consistent with respect to $\|\cdot\|_\infty$. By the Lax theorem, the scheme is convergent with respect to the $\|\cdot\|_\infty$. \qedhere
\end{enumerate} 
\end{enumerate} 
\end{proof}

Remark: 2+8+(4+6)


\newpage

\item (5 points) Refer to finite difference scheme in `Demonstration', by using a linear interpolation the price of the European call at $t=0$ is \$8.5889337. Meanwhile, the price derived from the Black-Scholes pricing formula is \$8.5916583. The percentage error is about 0.0317\%. Although the approximation sounds good, your boss wants to further reduce the percentage error. 

Give two suggestions, with different origins, to modify the current models (the mathematical model and the computational model/script) to your boss.

Mathematical proof nor programing are not required. However, you should explain your suggestions and state the possible tradeoffs if applicable.

The script of 'Demonstration 01' is shown below for your easy reference.

\begin{center}
\begin{Verbatim}[frame=single,fontsize=\footnotesize]
%Demonstration: 
%Using Matlab to Implement Explicit Finite Difference Scheme
%to Price a European Call

clear;

%basic parameters
S0=100; K=S0; r=0.05; sigma=0.15; T=1;
n=25; % number of interior spatial points
m=50; % number of temporal steps (no. of iterations)

%computed constants
Sr=S0*exp((r-sigma^2/2)*T+6*sigma*sqrt(T));
h=Sr/(n+1); % spatial step size
k=T/m; %temporal step size

%define the vector b
b=zeros(n,1);
b(end)=0.5*(sigma^2*n^2+r*n)*Sr;

%Define the iteration matrix G
G=diag(-(sigma^2*(1:n)^2+r)); % G=A11
G=G+diag(0.5*(sigma^2*(1:(n-1))+r).*(1:(n-1)),1); % G=A11+A12
G=G+diag(0.5*(sigma^2*(2:n)-r).*(2:n),-1); % G=A11+A12+A21 (=A)
G=eye(n,n)+k*G; % G=I+kA

%Iteration
v=max(h*(1:n)'-K,0) % initialize v, v=v0

for j=1:m
    v=G*v+k*b; % v_j=G*v_(j-1)+k*b
end
\end{Verbatim}
\end{center}

\begin{proof}
It is an open ended question so that any reasonable suggestions with explanations are fine.

Some possible modifications are given below.
\begin{itemize}
\item Increase the values of $n$ and/or $m$ if we believe the scheme is accurate of order $(2,1)$. It may be safe to keep the ratio of $k/h^2$ unchanged to reduce the chance that the scheme becomes unstable. Since the system of equations becomes larger and the number of iterations increases, the time consumed would be increased.

\item Modify the (artificial) boundary condition $V(S_r,T)=S_r$. Indeed, $S_r\approx 256$ which is not a good approximation of $V(S_r,T)=(S_r-100)^+\approx 156$ ($256/156\approx 1.64$). One may then change the boundary condition to $V(S_r,T)=(S_r-K)^+=S_r-K$. Alternatively, one may take a large value of $S_r$, for instance take $S_r=1000$ so that $V(S_r,T)=900$ ($1000/900\approx 1.11$). In the later approach, one may better increase $n$ (and hence $m$) if interpolation is to be used. \qedhere
\end{itemize}
\end{proof}

Remark: 3+3

\end{enumerate}

%\newpage

%\begin{center}
%{\it You may continue your solutions here.}
%\end{center}
\vfill

\begin{center}{\it The End}\end{center}

\end{document}
\textit{Solution.}
\begin{enumerate}
\item \begin{enumerate}
\item[(a)] \textbf{Example:} Consider the probability space $(\Omega, \mathcal{F}, \mathbb{P})$ where $\Omega = \{1,2,3,4\}$, $\mathcal{F} = 2^{\Omega}$, and $\mathbb{P}$ is any probability measure (for instance, uniform measure). Define a random variable $X:\Omega \to \mathbb{R}$ by
\[
X(\omega) = \omega.
\]
Now, define a function $f:\mathbb{R} \to \mathbb{R}$ as follows:
\[
f(x) = \begin{cases}
0, & x \in \{1, 2\}, \\
1, & x \in \{3, 4\}.
\end{cases}
\]
Then, the $\sigma$-algebra generated by $f(X)$, denoted $\sigma(f(X))$, is
\[
\sigma(f(X)) = \{\emptyset, \Omega, A, A^c\},
\]
where
\[
A = \{\omega \in \Omega : f(X(\omega)) = 1\} = \{3,4\}.
\]
The $\sigma$-algebra generated by $X$, $\sigma(X)$, is
\[
\sigma(X) = \{\emptyset, \Omega, \{\omega: X(\omega) \leq c\} \text{ for all } c \in \mathbb{R}\}.
\]
Specifically, it contains all subsets generated by the singleton sets $\{\omega\}$ because $X$ distinguishes all points. Since $A = \{3,4\}$ is a proper subset of $\Omega$, and $\sigma(f(X))$ only distinguishes between the sets where $f(X)$ is 0 or 1, it is strictly smaller than $\sigma(X)$. Moreover, since $\sigma(f(X))$ is generated by the partition $\{\{3,4\}, \{1,2\}\}$, it is non-trivial but strictly smaller than the full $\sigma$-algebra generated by $X$, which distinguishes all singleton points.

\item[(b)] \textbf{Answer:} No, the $\sigma$-algebra generated by $f(X)$ can never be strictly larger than $\sigma(X)$. This is because $f(X)$ is a measurable function of $X$, and the $\sigma$-algebra generated by $f(X)$, $\sigma(f(X))$, is always contained within $\sigma(X)$:
\[
\sigma(f(X)) \subseteq \sigma(X).
\]
Therefore, the $\sigma$-algebra generated by $f(X)$ cannot be strictly larger than that generated by $X$.
\end{enumerate}
\item \begin{enumerate}
\item[(a)] \textbf{Calculation of }$\mathbb{E}[\mathrm{Err}]$:

Recall that $\mathrm{Err} = Y - \mathbb{E}[Y|\mathcal{G}]$. Since $\mathbb{E}[Y|\mathcal{G}]$ is the conditional expectation of $Y$ given $\mathcal{G}$, it satisfies the property
\[
\mathbb{E}[\mathbb{E}[Y|\mathcal{G}]] = \mathbb{E}[Y].
\]
Using linearity of expectation:
\[
\mathbb{E}[\mathrm{Err}] = \mathbb{E}[Y - \mathbb{E}[Y|\mathcal{G}]] = \mathbb{E}[Y] - \mathbb{E}[\mathbb{E}[Y|\mathcal{G}]] = \mathbb{E}[Y] - \mathbb{E}[Y] = 0.
\]
\textbf{Answer:}
\[
\boxed{
\mathbb{E}[\mathrm{Err}] = 0.
}
\]

\item[(b)] \textbf{Proof of the variance decomposition:}

Given $X$ is $\mathcal{G}$-measurable, define $\mu := \mathbb{E}[Y - X]$.

We want to show:
\[
\mathrm{Var}(Y - X) = \mathbb{E}[\mathrm{Err}^2] + \mathbb{E}\left[(\mathbb{E}[Y|\mathcal{G}] - X - \mu)^2\right].
\]

Starting with the left side:
\[
\mathrm{Var}(Y - X) = \mathbb{E}[(Y - X - \mu)^2].
\]

Note that:
\[
Y - X - \mu = (Y - \mathbb{E}[Y|\mathcal{G}]) + (\mathbb{E}[Y|\mathcal{G}] - X - \mu) = \mathrm{Err} + (\mathbb{E}[Y|\mathcal{G}] - X - \mu).
\]

Since $\mathrm{Err}$ is orthogonal (uncorrelated) to any $\mathcal{G}$-measurable random variable (by properties of conditional expectation), we have:
\[
\mathbb{E}[\mathrm{Err} \cdot (\mathbb{E}[Y|\mathcal{G}] - X - \mu)] = 0.
\]

Therefore,
\[
\mathrm{Var}(Y - X) = \mathbb{E}\left[(\mathrm{Err} + (\mathbb{E}[Y|\mathcal{G}] - X - \mu))^2\right] = \mathbb{E}[\mathrm{Err}^2] + \mathbb{E}\left[(\mathbb{E}[Y|\mathcal{G}] - X - \mu)^2\right],
\]
as the cross term vanishes.

\textbf{Answer:}
\[
\boxed{
\mathrm{Var}(Y - X) = \mathbb{E}[\mathrm{Err}^2] + \mathbb{E}\left[(\mathbb{E}[Y|\mathcal{G}] - X - \mu)^2\right].
}
\]

\item[(c)] \textbf{Inequality for the variance of the error:}

Since the variance of the error is minimized when choosing the conditional expectation as the estimator, for any $\mathcal{G}$-measurable $X$, the following inequality holds:
\[
\mathrm{Var}(\mathrm{Err}) \leq \mathrm{Var}(Y - X).
\]

\textit{Proof:}

From part (b), setting $\mu = 0$ (since variance is unaffected by adding constants), we have:
\[
\mathrm{Var}(Y - X) = \mathbb{E}[\mathrm{Err}^2] + \mathbb{E}[(\mathbb{E}[Y|\mathcal{G}] - X)^2].
\]
Since the second term is non-negative,
\[
\mathrm{Var}(Y - X) \geq \mathbb{E}[\mathrm{Err}^2] = \mathrm{Var}(\mathrm{Err}),
\]
which completes the proof.

\textbf{Answer:}
\[
\boxed{
\mathrm{Var}(\mathrm{Err}) \leq \mathrm{Var}(Y - X).
}
\]
\end{enumerate}
\item
\begin{enumerate}
\item[(a)] \textbf{Limit of} 
\[
\sum_{j=0}^{n-1} (W_{t_{j+1}} - W_{t_j})(t_{j+1} - t_j) \quad \text{as} \quad \|\Pi\| \to 0.
\]
Since $\{W_{t_{j+1}} - W_{t_j}\}_{j=0}^{n-1}$ are mutually independent, and $W_{t_{j+1}} - W_{t_j}\sim N(0,t_{j+1}-t_j)$, we obtain
\begin{align*}
\mathbb{E}\left[\left(\sum_{j=0}^{n-1} (W_{t_{j+1}} - W_{t_j})(t_{j+1} - t_j) \right)^2\right]
&=\sum_{j=0}^{n-1}\mathbb{E}\left[\left(W_{t_{j+1}} - W_{t_j}\right)^2\right](t_{j+1} - t_j)\\
&=\sum_{j=0}^{n-1}(t_{j+1} - t_j)^2\\
&\leq\|\Pi\|\sum_{j=0}^{n-1}(t_{j+1} - t_j)=\|\Pi\|T\to 0 \text{ as } \|\Pi\|\to 0.
\end{align*}
\textbf{Solution:} 
\[
\boxed{
\lim_{\|\Pi\| \to 0} \sum_{j=0}^{n-1} (W_{t_{j+1}} - W_{t_j})(t_{j+1} - t_j) = 0 \quad \text{in } L^2(\Omega).
}
\]

\item[(b)] \textbf{Limit of} 
\[
\sum_{j=0}^{n-1} (t_{j+1} - t_j)^2 \quad \text{as} \quad \|\Pi\| \to 0.
\]
\textbf{Solution:} 
This sum depends only on the partition of the interval, independent of the Wiener process:
\[
\sum_{j=0}^{n-1} (t_{j+1} - t_j)^2\leq\|\Pi\|\sum_{j=0}^{n-1}(t_{j+1} - t_j)=\|\Pi\|T\to 0 \text{ as } \|\Pi\|\to 0.
\]
\[
\boxed{
\lim_{\|\Pi\| \to 0} \sum_{j=0}^{n-1} (t_{j+1} - t_j)^2 = 0.
}
\]

\item[(c)] \textbf{Limit of} 
\[
\sum_{j=0}^{n-1} (W_{t_{j+1}} - W_{t_j})^3 \quad \text{as} \quad \|\Pi\| \to 0.
\]
\textbf{Solution:} 
Since \(W_{t_{j+1}} - W_{t_j}\) are independent normal increments with mean zero, and
\[
(W_{t_{j+1}} - W_{t_j})^3\sim N(0, 15(t_{j+1}-t_j)^3)
\]
\begin{align*}
\mathbb{E}\left[\left(\sum_{j=0}^{n-1} (W_{t_{j+1}} - W_{t_j})^3\right)^2\right]&=
\sum_{j=0}^{n-1}\mathbb{E}\left[ (W_{t_{j+1}} - W_{t_j})^6\right]\\
&=15\sum_{j=0}^{n-1}(t_{j+1}-t_j)^3\\
&\leq 15\|\Pi\|^2T\to 0\text{ as }\|\Pi\|\to 0.
\end{align*}
Hence,
\[
\boxed{
\lim_{\|\Pi\| \to 0} \sum_{j=0}^{n-1} (W_{t_{j+1}} - W_{t_j})^3 = 0 \quad \text{in } L^2(\Omega).
}
\]
\end{enumerate}
\item \begin{enumerate}
        \item Change the order of integration, we obtain
        \begin{align*}        
                \mathbb{E}\bigg[ \int_0^T |t^2 W_t|^2 \,dt \bigg] 
                &= \int_0^T t^4 \mathbb{E}\bigg[ W_t^2 \bigg] \,dt 
                = \int_0^T t^4 \times t \,dt \\
                &= T^6/6 < \infty. 
        \end{align*}
        Here, we have used the property $W_t\sim N(0,t)$.
        \item \begin{equation*}
            \mathbb{E}[\xi(T)] = \mathbb{E}\bigg[ \int_0^T t^2W_t \,dW_t \bigg] = 0. 
        \end{equation*}
        
        \begin{equation*}
            \begin{split}
                \operatorname{Var}(\xi(T)) &=  \mathbb{E}[\xi(T)^2] - ( \mathbb{E}[\xi(T)])^2 \\
                &=  \mathbb{E}\bigg[ \big( \int_0^T t^2W_t \,dW_t \big)^2 \bigg] - 0 \\
                &= T^6/6. 
            \end{split}
        \end{equation*}
        
        \item Yes. 
    $$ X_t = X_0\exp( (\mu_1 - \frac{\sigma_1^2}{2})t + \sigma_1 W_t), $$
    $$ Y_t = Y_0\exp( (\mu_2 - \frac{\sigma_2^2}{2})t + \sigma_2 W_t). $$
    Then 
    $$ Z_t = Z_0\exp( (\mu_1 - \frac{\sigma_1^2}{2} + \mu_2 - \frac{\sigma_2^2}{2})t + (\sigma_1 + \sigma_2) W_t) $$
    with $Z_0 = X_0Y_0$. Hence, $Z_t$ is still a geometric Brownian motion. 
        
        \item Yes. 
    $$ U_t = \ln(X_t) - \ln(Y_t) $$
    Hence, 
    $$ dU_t = d\ln(X_t) - d\ln(Y_t) = (\mu_1 - \mu_2 + \frac{\sigma_2^2 - \sigma_1^2}{2})\,dt + (\sigma_1 - \sigma_2)\,dW_t.  $$
    \end{enumerate}
\end{enumerate}
\end{document}